\chapter{Distressed, Sovereign and Bank-Capital Credit}\label{ch:m2:distressed-sovereign-and-bank-capital-credit}

On Sunday 19 March 2023 the Swiss authorities announced that UBS would take over Credit
Suisse. The same day the supervisor, FINMA, said that the extraordinary government
support the deal required triggered a complete write-down of all of Credit Suisse's
Additional Tier 1 bonds, about CHF~16 billion, to zero. Credit Suisse's shareholders,
who rank below those bonds in a liquidation, received UBS shares worth CHF~3 billion.
The next day the European supervisors and resolution authority stated that in the
European Union equity absorbs losses first. In October 2025 a Swiss court held the
write-down unlawful; FINMA appealed. This chapter is about credit after things go
wrong: \omterm{def:m2:distressed-sovereign-and-bank-capital-credit:distressed}{distressed debt} and \omterm{def:m2:distressed-sovereign-and-bank-capital-credit:distressed}{exchange offers}, sovereigns that cannot be taken to a
bankruptcy court, and the bank bonds designed to absorb losses before the bank fails.

\section{Restructuring and exchange offers}

\begin{definition}[Distressed debt, exchange offer]\label{def:m2:distressed-sovereign-and-bank-capital-credit:distressed}
\emph{Distressed debt}\index{distressed debt} is debt of a borrower in or near default,
trading at a price that reflects an expected loss rather than a spread over the
risk-free rate. An \emph{exchange offer}\index{exchange offer} is a borrower's offer to
its creditors to swap their claims for new instruments, typically worth less, with
lower face value, lower coupons or longer maturities, sometimes with cash.
\end{definition}

A company in trouble can file for bankruptcy, where a court imposes a plan on all
creditors by class; or it can negotiate out of court, and \omterm{def:m2:distressed-sovereign-and-bank-capital-credit:distressed}{exchange offers} are how.
Out of court can be cheaper and faster, but each creditor decides for itself, and a
creditor who does not accept keeps its old claim. The value of an offer is the present
value of what it gives, at the yield at which the new instruments will trade after the
exchange, the \emph{exit yield}, compared with the old claim: the loss in present value
is the offer's haircut, often much larger than the cut in face value.

\begin{example}[An exchange offer]\label{ex:m2:distressed-sovereign-and-bank-capital-credit:offer}
A borrower offers, per 100 of old bonds, 50 of new 15-year bonds paying 4\% and 5 in
cash for those who tender. At an exit yield of 9\% the new bonds are worth 29.85 and
the package 34.85: a face-value haircut of 50\% but a present-value haircut of 65.2\%.
At an exit yield of 6\% the package is worth 45.29, at 12\% only 27.76
(\cref{fig:m2:distressed-sovereign-and-bank-capital-credit:offer}).
\end{example}

\begin{omfigure}
\begin{tikzpicture}
\begin{axis}[width=11cm,height=4.8cm,
  xlabel={exit yield (\%)},ylabel={value per 100 of old face},
  tick label style={font=\small},label style={font=\small},
  xmin=5,xmax=15,ymin=15,ymax=60,grid=major,grid style={gray!25},
  legend columns=2,legend style={font=\footnotesize,at={(0.5,-0.3)},anchor=north,draw=none,fill=none,/tikz/every even column/.append style={column sep=8pt}},
  table/col sep=comma]
\addplot[omDef,very thick] table[x=yield,y=withcash]{figdata/markets-2/26-distressed-sovereign-and-bank-capital-credit/offer.csv};
\addplot[omThm,very thick,dashed] table[x=yield,y=bonds]{figdata/markets-2/26-distressed-sovereign-and-bank-capital-credit/offer.csv};
\addplot[only marks,mark=*,mark size=1.6pt,black,forget plot] coordinates {(6,45.29) (9,34.85) (12,27.76)};
\legend{new bonds and cash,new bonds alone}
\end{axis}
\end{tikzpicture}

\omcaption{Value of the chapter's \omterm{def:m2:distressed-sovereign-and-bank-capital-credit:distressed}{exchange offer}, per 100 of old face, against the exit
yield at which the new 15-year 4\% bonds will trade: with the 5 of cash paid to those
who tender, and without it, which is what a holder bound by a \omterm{def:m2:distressed-sovereign-and-bank-capital-credit:cac}{collective action clause}
receives. The dots are \cref{ex:m2:distressed-sovereign-and-bank-capital-credit:offer}.
Illustrative; data: the chapter's tutorial.}\label{fig:m2:distressed-sovereign-and-bank-capital-credit:offer}
\end{omfigure}

\section{Sovereign debt: collective action clauses and holdouts}

A sovereign has no bankruptcy court. Its bonds can be restructured only by agreement,
and every creditor who does not agree can refuse, keep its claim and sue.

\begin{definition}[Collective action clause, holdout creditor]\label{def:m2:distressed-sovereign-and-bank-capital-credit:cac}
A \emph{collective action clause}\index{collective action clause} (CAC) in a bond lets a
qualified majority of its holders change its payment terms for all holders, the
minority included. A \emph{holdout creditor}\index{holdout creditor} is one that does
not accept a restructuring and keeps, or litigates on, its original claim.
\end{definition}

The IMF describes the traditional clause as series by series: typically 75\% of the
holders of one bond can amend it, so a creditor who buys a blocking stake in one series
can hold that series out. Newer clauses aggregate: in the model clauses of the
International Capital Market Association, a single vote of 75\% of the principal of all
the series affected binds every series, provided all are offered the same terms, or,
in a two-limb version, 66\,2/3\% in aggregate and more than half of each series.

\begin{definition}[Pari passu clause]\label{def:m2:distressed-sovereign-and-bank-capital-credit:pp}
A \emph{pari passu clause}\index{pari passu clause} states that a bond ranks equally
with the issuer's other unsecured debt of the same kind.
\end{definition}

Two cases shaped current practice. Argentina defaulted in 2001 and restructured about
93\% of its external debt in exchanges in 2005 and 2010. Holdouts led by NML Capital,
with claims of about USD~1.6 billion, sued in New York and won a reading of the \omterm{def:m2:distressed-sovereign-and-bank-capital-credit:pp}{pari
passu clause} as requiring ratable payment: Argentina could not pay the exchanged bonds
without paying the holdouts in full at the same time, and the banks in the payment chain
were barred from passing on payments. The Supreme Court declined to hear Argentina's
appeal in June 2014, and at the end of July 2014 the exchanged bondholders, who had
settled for about 27 cents on the dollar in present value, went unpaid. Greece in 2012
did the opposite: a law passed in February introduced \omterm{def:m2:distressed-sovereign-and-bank-capital-credit:cac}{collective action clauses} into
its Greek-law bonds, binding if two thirds of the participating holders agreed. Holders
of about 85\% of the EUR~206 billion of privately held bonds accepted, the clauses
made the exchange binding on the Greek-law bonds, and new bonds of about EUR~100 billion
of face value replaced the old. About EUR~6 billion of English-law bonds held out; in
May 2012 Greece repaid a maturing EUR~435 million English-law bond in full.

Why hold out? A holdout's value is a probability tree
(\cref{fig:m2:distressed-sovereign-and-bank-capital-credit:tree}): the sovereign may
pay the holdouts in full to be rid of them, which is likelier the fewer they are, as
Greece did; if not, a lawsuit may win the whole claim years later, as NML won its case; or
the bond may stay in default. The more others tender, the better holding out looks,
which is the free-rider problem CACs exist to solve.

\begin{omfigure}
\centering
\begin{tikzpicture}[font=\small,
  nd/.style={draw,thick,rounded corners=2pt,align=center,inner sep=3pt,minimum height=0.8cm},
  arr/.style={-{Stealth[length=2.2mm]},thick}]
\node[nd,draw=omDef,fill=omDef!8] (h) at (0,0) {hold out\\(participation $p$)};
\node[nd,draw=omThm,fill=omThm!8] (paid) at (5,1.4) {paid in full next year\\worth 91.74};
\node[nd,draw=gray,fill=gray!8] (no) at (5,-1.2) {not paid soon};
\node[nd,draw=omProp,fill=omProp!8] (law) at (10.4,-0.3) {lawsuit recovers 130\\in six years: 77.51};
\node[nd,draw=gray,fill=gray!15] (stuck) at (10.4,-2.1) {stays in default\\worth 15};
\draw[arr] (h) -- node[above,sloped,font=\footnotesize] {$p^8$} (paid);
\draw[arr] (h) -- node[below,sloped,font=\footnotesize] {$1 - p^8$} (no);
\draw[arr] (no) -- node[above,sloped,font=\footnotesize] {0.3} (law);
\draw[arr] (no) -- node[below,sloped,font=\footnotesize] {0.7} (stuck);
\end{tikzpicture}

\omcaption{The holdout's probability tree in the chapter's example, values per 100 of old
face discounted at 9\%. The chance of being paid in full soon rises steeply with the
share of creditors who tendered, since the sovereign can afford to pay a small
remainder. Illustrative.}\label{fig:m2:distressed-sovereign-and-bank-capital-credit:tree}
\end{omfigure}

\begin{proposition}[The free-rider band]\label{prop:m2:distressed-sovereign-and-bank-capital-credit:band}
Let $T$ be the value of tendering, $H(p)$ the value of holding out when a share $p$
tenders, increasing in $p$, and $p^*$ the participation at which $H(p^*) = T$. Without a
\omterm{def:m2:distressed-sovereign-and-bank-capital-credit:cac}{collective action clause}, holding out pays whenever $p > p^*$. With an aggregated clause
of threshold $\tau > p^*$, a holdout is bound to the offer once $p \geq \tau$ and receives
only what the offer gives the bound; holding out pays only for $p^* < p < \tau$, and not at
all if $\tau \leq p^*$.
\end{proposition}

\begin{proof}
Below $\tau$ the holdout keeps its claim, worth $H(p)$, which exceeds $T$ exactly when
$p > p^*$. At or above $\tau$ the clause amends its bond to the offer's terms, which are
worth at most $T$ and less by any cash reserved for those who tendered.
\end{proof}

\begin{example}[Where holding out pays]\label{ex:m2:distressed-sovereign-and-bank-capital-credit:band}
In the tree of \cref{fig:m2:distressed-sovereign-and-bank-capital-credit:tree}, holding
out is worth 33.75 if no one tenders, 34.73 at 60\% participation and 38.97 at 74\%,
against 34.85 for tendering. It pays from a participation of 60.9\%. With a 75\%
aggregated clause, at 75\% or more the holdout is bound and receives the new bonds
alone, 29.85: the band in which holding out pays is 60.9\% to 75\%. Without a clause
it pays at any participation above 60.9\%, and at 90\% it is worth 58.72
(\cref{fig:m2:distressed-sovereign-and-bank-capital-credit:holdout}).
\end{example}

\begin{omfigure}
\begin{tikzpicture}
\begin{axis}[width=11cm,height=5cm,
  xlabel={participation in the exchange (\%)},ylabel={value per 100 of old face},
  tick label style={font=\small},label style={font=\small},
  xmin=0,xmax=100,ymin=20,ymax=95,grid=major,grid style={gray!25},
  legend columns=3,legend style={font=\footnotesize,at={(0.5,-0.3)},anchor=north,draw=none,fill=none,/tikz/every even column/.append style={column sep=8pt}},
  table/col sep=comma]
\addplot[omThm,very thick,dashed] table[x=p,y=nocac]{figdata/markets-2/26-distressed-sovereign-and-bank-capital-credit/holdout.csv};
\addplot[omDef,very thick,const plot] table[x=p,y=cac]{figdata/markets-2/26-distressed-sovereign-and-bank-capital-credit/holdout.csv};
\addplot[black,thick] table[x=p,y=tender]{figdata/markets-2/26-distressed-sovereign-and-bank-capital-credit/holdout.csv};
\addplot[gray,thin,dotted,forget plot] coordinates {(60.9,20) (60.9,95)};
\addplot[gray,thin,dotted,forget plot] coordinates {(75,20) (75,95)};
\node[font=\footnotesize,anchor=north east] at (axis cs:60.4,93) {$p^* = 60.9\%$};
\node[font=\footnotesize,anchor=north west] at (axis cs:75.5,93) {CAC 75\%};
\legend{hold out: no CAC,hold out: 75\% CAC,tender}
\end{axis}
\end{tikzpicture}

\omcaption{Value of holding out against the share of creditors who tender, without a
\omterm{def:m2:distressed-sovereign-and-bank-capital-credit:cac}{collective action clause} and with a 75\% aggregated clause, and the value of tendering.
Holding out pays above 60.9\% participation; the clause removes the gain from 75\%,
where the holdout is bound to the offer and loses the cash paid to those who tendered.
Illustrative; data: the chapter's tutorial.}\label{fig:m2:distressed-sovereign-and-bank-capital-credit:holdout}
\end{omfigure}

\section{Bank capital: contingent convertibles}

A bank that fails in a crisis may be too important to liquidate. After 2008 the
regulators' answer was to make some of its creditors absorb losses while the bank
keeps operating.

\begin{definition}[Contingent convertible bond, additional tier 1]\label{def:m2:distressed-sovereign-and-bank-capital-credit:coco}
A \emph{contingent convertible bond}\index{contingent convertible bond} (CoCo) is a
bank bond that is converted into shares or written down when a trigger is hit, such as
the bank's common equity ratio falling below a level set in the bond. \emph{Additional
tier 1}\index{additional tier 1} (AT1) capital consists of perpetual instruments,
usually CoCos, whose coupons the bank may cancel and which absorb losses while the
bank is a going concern; they rank above common equity and below all other debt in a
liquidation.
\end{definition}

\begin{definition}[Point of non-viability, bail-in]\label{def:m2:distressed-sovereign-and-bank-capital-credit:pon}
The \emph{point of non-viability}\index{point of non-viability} is the moment an
authority decides that a bank would fail without a write-down of its capital
instruments or public support. A \emph{bail-in}\index{bail-in} is the write-down or
conversion into equity of a failing bank's liabilities by an authority, so that
creditors instead of taxpayers bear its losses.
\end{definition}

The Basel Committee required in January 2011 that, from 2013, all capital instruments
other than common equity issued by internationally active banks can be written off or
converted into common equity, at the authority's option, at the earlier of a decision
that a write-off is necessary to keep the bank viable and a decision to inject public
capital or equivalent support without which it would not be. An AT1 bond therefore
carries two triggers: its own ratio trigger, while the bank is a going concern, and
the authority's at the \omterm{def:m2:distressed-sovereign-and-bank-capital-credit:pon}{point of non-viability}. Its holders are paid a high coupon for
the risk that either is hit, and for the risk that the bank simply stops paying the
coupon, which it may do without defaulting.

\begin{omfigure}
\centering
\begin{tikzpicture}[font=\small]
\def\w{3.2}
\foreach \x/\t in {0/{order of the EU statement},6.4/{Credit Suisse, 2023}} \node[font=\small\bfseries] at (\x+\w/2,5.3) {\t};
\foreach \x in {0,6.4} {
  \fill[gray!12] (\x,3.2) rectangle (\x+\w,4.8); \node at (\x+\w/2,4.0) {senior bail-in debt};
  \fill[omProp!20] (\x,2.4) rectangle (\x+\w,3.2); \node at (\x+\w/2,2.8) {Tier 2};
  \fill[omThm!25] (\x,1.2) rectangle (\x+\w,2.4); \node at (\x+\w/2,1.8) {AT1};
  \fill[omDef!25] (\x,0) rectangle (\x+\w,1.2); \node at (\x+\w/2,0.6) {common equity};
  \draw[thick] (\x,0) rectangle (\x+\w,4.8);
  \foreach \y in {1.2,2.4,3.2} \draw[thick] (\x,\y) -- (\x+\w,\y);
}
\draw[-{Stealth[length=2.5mm]},very thick,omDef] (3.5,0.2) -- node[right,font=\footnotesize,align=left] {1. equity\\first} (3.5,1.1);
\draw[-{Stealth[length=2.5mm]},very thick,omThm] (3.5,1.3) -- node[right,font=\footnotesize,align=left] {2. then\\AT1} (3.5,2.3);
\draw[-{Stealth[length=2.5mm]},very thick,omThm] (9.9,1.3) -- node[right,font=\footnotesize,align=left] {AT1 written\\to zero} (9.9,2.3);
\node[font=\footnotesize,align=left,anchor=west] at (9.8,0.6) {shareholders\\paid CHF~3 bn};
\end{tikzpicture}

\omcaption{The loss-absorbing stack of a bank. The EU authorities stated in March 2023
that common equity absorbs losses first and AT1 only after its full use; in Credit
Suisse's case the AT1 bonds were written down under their contractual viability trigger
and an emergency ordinance, while shareholders received UBS shares. Schematic; heights
not to scale.}\label{fig:m2:distressed-sovereign-and-bank-capital-credit:stack}
\end{omfigure}

\section{The 2023 write-down}

Credit Suisse's AT1 bonds said that they would be written down completely in a
\emph{viability event}, in particular if the bank received extraordinary government
support. On 19 March 2023 it received liquidity assistance loans backed by a federal
default guarantee, and an emergency ordinance of the Federal Council authorised FINMA to
order the write-down of AT1 capital; FINMA ordered it, and the bonds' nominal value of
about CHF~16 billion became zero, raising the bank's core capital. Tier 2 bonds were not
written down. Under the merger terms, Credit Suisse shareholders received one UBS share
for every 22.48 of theirs, CHF~3 billion in all
(\cref{fig:m2:distressed-sovereign-and-bank-capital-credit:stack}).

To bondholders who read the capital stack as a liquidation order, this was a reversal:
the claim senior to equity received nothing while equity received something. The
European supervisors' statement of 20 March, that equity absorbs losses first and AT1
only after its full use, told investors that the European framework would not do the
same. The write-down followed the bonds' contractual terms rather than a liquidation
order, and whether those terms and the ordinance allowed it became a matter for the
courts: on 1 October 2025 the Federal Administrative Court, in the first of about 360
proceedings, held that the contractual viability event had not occurred and that the
write-down lacked a legal basis; the judgment was not final, and FINMA appealed to the
Federal Supreme Court.

\begin{dated}{2026-09}{The Credit Suisse AT1 litigation}\label{dat:m2:distressed-sovereign-and-bank-capital-credit:at1}
The Federal Administrative Court's judgment of 1 October 2025 annulled FINMA's order of
19 March 2023 writing down about CHF~16.5 billion of AT1 bonds; FINMA announced on 15
October 2025 that it would appeal to the Federal Supreme Court. No final ruling had been
found when this chapter was written.
\end{dated}

\begin{example}[Two orders of loss]\label{ex:m2:distressed-sovereign-and-bank-capital-credit:stack}
An illustrative bank has 45 of common equity, 16 of AT1, 12 of Tier 2 and 80 of senior
debt eligible for \omterm{def:m2:distressed-sovereign-and-bank-capital-credit:pon}{bail-in}, in billions. A loss of 20 absorbed in the order of the EU
statement leaves equity with 25 and every bond intact. An AT1 write-down under a
contractual trigger removes all 16 of the AT1 and leaves the equity with 45.
\end{example}

\section{Tutorial: an exchange offer against holding out}\label{tut:m2:distressed-sovereign-and-bank-capital-credit:recovery}

\begin{tutorial}
\textbf{Goal.} Value an \omterm{def:m2:distressed-sovereign-and-bank-capital-credit:distressed}{exchange offer} at exit yields, value holding out with a
probability tree, find the participation at which holding out starts to pay, and see
what an aggregated \omterm{def:m2:distressed-sovereign-and-bank-capital-credit:cac}{collective action clause} changes.
\textbf{End state:} \cref{fig:m2:distressed-sovereign-and-bank-capital-credit:offer,fig:m2:distressed-sovereign-and-bank-capital-credit:holdout},
\cref{ex:m2:distressed-sovereign-and-bank-capital-credit:offer,ex:m2:distressed-sovereign-and-bank-capital-credit:band,ex:m2:distressed-sovereign-and-bank-capital-credit:stack}
and the numbers of the weekend problem.
\begin{tutsteps}
\item \textbf{The offer}: bond values at an exit yield, haircuts, the package.
\omcode{code/firm/recovery/firm_recovery.py}{13}{30}{Bond values, haircuts and the exchange package.}\label{lst:m2:distressed-sovereign-and-bank-capital-credit:offer}
\item \textbf{The holdout}: the probability tree, the clause, the break-even participation
and the loss-absorbing stack.
\omcode{code/firm/recovery/firm_recovery.py}{33}{69}{The holdout's tree, the CAC, the break-even and loss absorption.}\label{lst:m2:distressed-sovereign-and-bank-capital-credit:holdout}
\item \textbf{Run} \texttt{restructuring\_demo.tutorial()},
\texttt{restructuring\_demo.capital()} and \texttt{fig\_restructuring.py}.
\end{tutsteps}
\textbf{What to change next.} Replace the aggregated clause by series-by-series
clauses on four bonds, and let a holdout buy 26\% of one series: find the participation
the sovereign now needs. Then make the chance of being paid soon depend on the
sovereign's cash rather than on participation alone.
\end{tutorial}

\section{Build: restructuring-offer valuation}\label{bld:m2:distressed-sovereign-and-bank-capital-credit:recovery}

\begin{build}
\bfield{Purpose} The miniature firm's distressed desk decides whether to tender or hold
out, and its credit desk values bank capital bonds under both orders of loss.
\bfield{Interface} \texttt{bond\_pv(face, coupon, years, y)}; \texttt{npv\_haircut(value,
claim)}; \texttt{Offer(new\_face, coupon, years, cash).value(y)};
\texttt{Holdout(paid\_soon, lawsuit\_prob, lawsuit\_pv, stuck, power).value(p)};
\texttt{holdout\_payoff(p, holdout, bound\_value, cac)};
\texttt{breakeven\_participation(holdout, tender\_value)}; \texttt{absorb(stack, loss)}.
\bfield{Rules} Annual coupons; values per 100 of old face at one exit yield; the chance
of being paid soon is participation to a power; a CAC binds all holders at or above
its threshold, without the tender cash; losses absorbed in the order of the stack.
\bfield{Acceptance tests} \texttt{code/firm/recovery/tests/}: a bond at its coupon yield
is worth par; the offer's value falls with the exit yield; the holdout's value rises
with participation and the clause binds; the break-even solves $H(p) = T$; the stack
absorbs in order.
\bfield{Stretch} Series-by-series clauses and blocking stakes; a game between creditors
with a distribution of beliefs about participation; recovery analysis of a company's
capital structure in a bankruptcy plan; AT1 pricing with coupon cancellation and
extension risk (One Quant Book 6).
\end{build}

\begin{omsources}
\item FINMA, press releases of 19 and 23 March 2023 and of 15 October 2025; Swiss
Federal Administrative Court, media release of 14 October 2025.
\item UBS, media release of 19 March 2023.
\item SRB, EBA and ECB Banking Supervision, statement of 20 March 2023.
\item Basel Committee on Banking Supervision, press release of 13 January 2011 on loss
absorbency at the point of non-viability.
\item BIS Quarterly Review, December 2012, box on governing law and the Greek debt
restructuring.
\item IMF, ``Strengthening the contractual framework to address collective action
problems in sovereign debt restructuring'', October 2014.
\end{omsources}

\section{Exercises}

\begin{exercise}[$\star$]\label{exo:m2:distressed-sovereign-and-bank-capital-credit:1}
Per 100 of old bonds, an offer gives 60 of new 10-year 5\% bonds. What is its face
haircut, and its present-value haircut at an exit yield of 8\%?
\end{exercise}

\begin{exercise}[$\star$]\label{exo:m2:distressed-sovereign-and-bank-capital-credit:2}
What is the difference between a series-by-series and a single-limb aggregated
\omterm{def:m2:distressed-sovereign-and-bank-capital-credit:cac}{collective action clause}?
\end{exercise}

\begin{exercise}[$\star$]\label{exo:m2:distressed-sovereign-and-bank-capital-credit:3}
Name the two triggers an AT1 bond carries, and what each depends on.
\end{exercise}

\begin{exercise}[$\star\star$]\label{exo:m2:distressed-sovereign-and-bank-capital-credit:4}
With the Greek threshold of two thirds instead of 75\%, over what range of participation
does holding out pay in the chapter's example?
\end{exercise}

\begin{exercise}[$\star\star$]\label{exo:m2:distressed-sovereign-and-bank-capital-credit:5}
Why did the ratable-payment reading of the \omterm{def:m2:distressed-sovereign-and-bank-capital-credit:pp}{pari passu clause} make Argentina's
restructuring harder, and how have newer bond contracts responded?
\end{exercise}

\begin{exercise}[$\star\star$]\label{exo:m2:distressed-sovereign-and-bank-capital-credit:6}
Why did Greece repay its English-law holdouts in full in 2012 when the Greek-law
bondholders took a large loss?
\end{exercise}

\begin{exercise}[$\star\star\star$]\label{exo:m2:distressed-sovereign-and-bank-capital-credit:7}
\emph{Coding.} Raise the lawsuit's probability of success from 0.3 to 0.5 and find the
new break-even participation. What does a stronger legal position do to the free-rider
band?
\end{exercise}

\begin{exercise}[$\star\star\star$]\label{exo:m2:distressed-sovereign-and-bank-capital-credit:8}
\emph{Find the flaw.} ``AT1 bonds rank above shares, so AT1 holders cannot lose
everything while shareholders keep something.'' Correct it.
\end{exercise}

\section{Problem: Hold Out or Tender}

\begin{problem}[{Weekend problem --- the free-rider band}]\label{pb:m2:distressed-sovereign-and-bank-capital-credit:1}
A sovereign in default offers, per 100 of old bonds, 50 of new 15-year 4\% bonds and 5
in cash for those who tender. The new bonds are expected to trade at a 9\% yield. A fund
holding 100 million of old bonds weighs holding out: the sovereign may pay remaining
holdouts in full next year, with a probability equal to the participation raised to the
eighth power; if not, a lawsuit has a 30\% chance of recovering 130 per 100 in six
years; otherwise the bonds stay in default, worth 15. The bonds carry a single-limb
aggregated \omterm{def:m2:distressed-sovereign-and-bank-capital-credit:cac}{collective action clause} with a 75\% threshold.

\medskip\noindent\textbf{Part I --- The offer.}
\begin{enumerate}
\item What is the offer worth per 100, with and without the cash?
\item What are its face and present-value haircuts?
\item What is it worth at exit yields of 6\% and 12\%?
\item Why does the exit yield matter so much?
\item What does the fund receive if it tenders its 100 million?
\end{enumerate}

\medskip\noindent\textbf{Part II --- Holding out.}
\begin{enumerate}[resume]
\item Value each branch of the tree.
\item What is holding out worth if no one tenders, and at 60\% and 74\% participation?
\item At what participation does holding out start to pay?
\item What happens to a holdout at 75\% participation or more?
\item What would holding out be worth at 90\% participation without a clause?
\end{enumerate}

\medskip\noindent\textbf{Part III --- Strategy.}
\begin{enumerate}[resume]
\item If every creditor reasoned like the fund, could participation settle between 60.9\%
and 75\%?
\item Why does the cash for tendering matter once the clause binds?
\item How could the fund block the clause, and why is that harder with aggregation?
\item What did Greece's English-law holdouts receive, and why?
\item What did Argentina's exchanged bondholders suffer from the holdouts' victory?
\end{enumerate}

\medskip\noindent\textbf{Part IV --- Judgement.}
\begin{enumerate}[resume]
\item Why is the probability of being paid soon hard to estimate?
\item Should a fund that buys defaulted bonds cheaply to litigate be allowed to block a
restructuring?
\item How would you hedge the fund's position while it waits?
\item State the \emph{named result}: the participation band in which holding out pays,
and what the clause does above it.
\item In one sentence: what problem do \omterm{def:m2:distressed-sovereign-and-bank-capital-credit:cac}{collective action clauses} solve?
\end{enumerate}
\end{problem}

\section{Interview questions}

\begin{interviewq}[$\star$ \iqroles{trader, researcher}]\label{iq:m2:distressed-sovereign-and-bank-capital-credit:1}
How is a sovereign debt restructuring different from a corporate bankruptcy?
\end{interviewq}

\begin{interviewq}[$\star$ \iqroles{trader, bank}]\label{iq:m2:distressed-sovereign-and-bank-capital-credit:2}
What is an AT1 bond, and what are its risks?
\end{interviewq}

\begin{interviewq}[$\star\star$ \iqroles{researcher}]\label{iq:m2:distressed-sovereign-and-bank-capital-credit:3}
How would you value a distressed \omterm{def:m2:distressed-sovereign-and-bank-capital-credit:distressed}{exchange offer}?
\end{interviewq}

\begin{interviewq}[$\star\star$ \iqroles{trader, researcher}]\label{iq:m2:distressed-sovereign-and-bank-capital-credit:4}
Explain the holdout problem and how \omterm{def:m2:distressed-sovereign-and-bank-capital-credit:cac}{collective action clauses} address it.
\end{interviewq}

\begin{interviewq}[$\star\star$ \iqroles{risk, bank}]\label{iq:m2:distressed-sovereign-and-bank-capital-credit:5}
What did the Credit Suisse AT1 write-down teach investors about bank capital?
\end{interviewq}

\begin{interviewq}[$\star\star\star$ \iqroles{developer}]\label{iq:m2:distressed-sovereign-and-bank-capital-credit:6}
Design a tool that tracks, for a sovereign's outstanding bonds, the governing law, the
type of \omterm{def:m2:distressed-sovereign-and-bank-capital-credit:cac}{collective action clause} and the votes needed to bind each series.
\end{interviewq}
